\section{Exercícios integrados e projetos} \label{sec:ExercicioIntetradosProjetos}


\subsection*{Nível 1 -- Estruturas Numéricas Básicas}

\begin{enumerate}[label=\thesection.\arabic*]

\item Seja o vetor\footnote{O arquivo .csv com o conteúdo deste vetor está disponível 
em \href{https://github.com/JonathaCosta-IA/NM/blob/main/B-NM_codes/SolListaExerciciosPropostos/z_AplicMN_eng/dados_temp.csv}{Vetor}} de medições de temperatura(°C) a seguir.
\[
t(^{\circ}\mathrm{C}) = \begin{bmatrix}
	37.454011884736246 \\
	95.071430640991622 \\
	73.199394181140505 \\
	59.865848419703660 \\
	15.601864044243651 \\
	15.599452033620265 \\
	5.808361216819947 \\
	86.617614577493512 \\
	60.111501174320878 \\
	70.807257779604555 \\
	2.058449429580245 \\
	96.990985216199434 \\
	83.244264080042171 \\
	21.233911067827616 \\
	18.182496720710063 \\
	18.340450985343381 \\
	30.424224295953771 \\
	52.475643163223786 \\
	43.194501864211574 \\
	29.122914019804192 \\
	61.185289472237947 \\
	13.949386065204184 \\
	29.214464853521815 \\
	36.636184329369172 \\
	45.606998421703594
\end{bmatrix}
\]
Este vetor contém 16 dígitos que representam registros térmicos de um forno industrial. 
Considerando que os dados deverão ser armazenados em precisão simples, realize as seguintes tarefas:
	\begin{itemize}
	\item Converta todos os elementos para precisão simples (float com 8 dígitos).
	\item Calcule o erro absoluto e relativo médio.
	\item Analise o impacto acumulado em somatórios.
	\end{itemize}
\item Considere os dados da questão anterior.

\begin{itemize}
    \item Calcule o polinômio de interpolação $p(x)$ utilizando Lagrange.
    \item Calcule o polinômio de interpolação $p(x)$ utilizando Newton.   
    \item Calcule o polinômio de interpolação $p(x)$ utilizando regressão linear.
    \item Calcule o polinômio de interpolação $p(x)$ utilizando regressão polinomial com graus(2,3,4).
    \item Compare erro e custo computacional para um ponto de teste $t=50^{\circ}\mathrm{C}$.
\end{itemize}
	
\item 
\end{enumerate}

\subsection*{Nível 3 -- Métodos Numéricos Aplicados à Irradiação Solar}

\begin{enumerate}[label=\thesection.\arabic*]

\item Seja a matriz\footnote{O arquivo .csv com o conteúdo desta matriz está disponível 
em \href{https://github.com/JonathaCosta-IA/NM/blob/main/B-NM_codes/SolListaExerciciosPropostos/z_AplicMN_eng/dados_radiacao.csv}{MatrizRadiacao}} de perfis horários médios mensais de irradiação solar global \((\mathrm{Wh/m^2})\) apresentada a seguir:

\[
\small
A =
\begin{bmatrix}
134 & 114 & 127 & \cdots & 177 \\
245 & 238 & 259 & \cdots & 266 \\
324 & 327 & 337 & \cdots & 368 \\
\vdots & \vdots & \vdots & \ddots & \vdots \\
310 & 282 & 277 & \cdots & 331 \\
93 & 95 & 74 & \cdots & 64
\end{bmatrix}
\]

Cada linha representa um intervalo horário entre \(6\mathrm{h}\) e \(18\mathrm{h}\), enquanto cada coluna representa um mês do ano.

Este conjunto de dados representa medições utilizadas em sistemas de monitoramento e previsão de geração fotovoltaica. 
Considerando aplicações em métodos numéricos, realize as seguintes tarefas:

\begin{itemize}

\item Converta todos os elementos da matriz para precisão simples (\texttt{float} com 8 dígitos).

\item Calcule o erro absoluto médio e o erro relativo médio da conversão.

\item Analise o impacto da redução de precisão no cálculo da energia solar acumulada ao longo do dia.

\end{itemize}

%

\item Considere os dados da questão anterior referentes ao mês de agosto.

\begin{itemize}

\item Calcule o polinômio de interpolação \(p(t)\) utilizando o método de Lagrange.

\item Calcule o polinômio de interpolação \(p(t)\) utilizando o método de Newton.

\item Calcule o modelo aproximado utilizando regressão linear.

\item Calcule modelos de regressão polinomial de graus \(2\), \(3\) e \(4\).

\item Compare erro numérico e custo computacional dos métodos para estimar a irradiação solar no instante:

\[
t = 10.5\mathrm{h}
\]

\end{itemize}

%

\item Considere os dados referentes ao mês de outubro.

A energia solar diária acumulada pode ser aproximada por:

\[
E = \int_{6}^{18} I(t)\,dt
\]

em que: \(I(t)\) representa a irradiação solar horária.

\begin{itemize}

\item Calcule a integral numérica utilizando a Regra do Trapézio.

\item Calcule a integral numérica utilizando a Regra de Simpson \(1/3\).

\item Compare os resultados obtidos e discuta a precisão numérica dos métodos.

\end{itemize}

%

\item Considere os dados do mês de julho.

\begin{itemize}

\item Aproximar numericamente a derivada da irradiação solar no instante \(t=12\mathrm{h}\) utilizando:

\begin{itemize}

\item diferença progressiva;

\item diferença regressiva;

\item diferença central.

\end{itemize}

\item Interpretar fisicamente o significado da taxa de variação obtida.

\item Comparar a estabilidade numérica das aproximações calculadas.

\end{itemize}

%
\item Considere a matriz de covariância associada aos dados mensais da matriz \(A\).

\begin{itemize}

\item Determine numericamente os autovalores e autovetores principais.

\item Identifique a componente principal responsável pela maior variabilidade dos dados.

\item Discuta aplicações da decomposição espectral em:

\begin{itemize}

\item previsão de geração fotovoltaica;

\item compressão de dados;

\item sistemas inteligentes de monitoramento energético.

\end{itemize}

\end{itemize}

%

\item Considere que a potência instantânea gerada por um painel fotovoltaico é dada por:

\[
P(t)=\eta A_p I(t)
\]

em que:

\begin{itemize}

\item \(\eta\) representa a eficiência do painel fotovoltaico;

\item \(A_p\) representa a área do módulo fotovoltaico;

\item \(I(t)\) representa a irradiação solar incidente.

\end{itemize}

Utilizando métodos numéricos, realize as seguintes tarefas:

\begin{itemize}

\item Estime a energia total gerada em um dia típico.

\item Determine numericamente o ponto de máxima geração.

\item Utilize spline cúbica para prever valores de geração em horários não amostrados.

\item Compare spline cúbica e interpolação polinomial global quanto à estabilidade numérica.

\end{itemize}

\end{enumerate}